% Requires mathrsfs for \mathscr{P} and TikZ (arrows.meta) for Figure~\ref{fig:reverse-fold}.
\subsection{Further results}\label{sec:additional-results}
In \S\ref{sec:overview}, we explained how folds of a reference structure
prove our main theorem. We now use this approach to establish further
exponential succinctness gaps: from $\FOthree$ to $\FOtwo$, from
$\CoR$ to $\FOtwo$, and from three-attribute relational algebra to
matrix query languages. To explain these extensions, we first summarize
the counting argument of \S\ref{sec:overview}.

Recall that the folds in \S\ref{sec:overview} are indexed by colors $\lambda\in\Lambda=2^{[k]}$.
For a circuit $C \equiv G_k$,
let $D_g\subseteq\Lambda$ be the set of colors whose folds are noticed by gate $g$.
Every color belongs to some $D_g$: otherwise, the circuit
could not detect the change at the read edge left unchanged by that
fold (\S\ref{ov:types}). Thus $\bigcup_{g\in C}D_g=\Lambda$, and
counting gives
\begin{equation}\tag{$\dagger$}\label{eq:fold-counting}
  2^k\le \sum_{g\in C}|D_g|.
\end{equation}
For \CoR, Boolean operations and converse notice no folds, while each
composition notices at most two (\S\ref{ov:twocolors}). Hence
$\sum_{g\in C}|D_g|\le 2\gcost C$, recovering Theorem~\ref{thm:main}.
The applications below use the same counting argument, with different
bounds on the folds noticed by the target operations and, when needed,
a different reference structure.

\paragraph{\texorpdfstring{{\boldmath\FOthree\ to \FOtwo}}{FO3 to FO2}.}
The same query $G_k$ is also \FOtwo-definable: the DNF expansion in
\S\ref{sec:main-results} gives
\[
 G_k(x,y)\equiv\bigvee_{A\subseteq[k]}
 \bigl[H_A(x,y)\land\exists y\,H_{A^c}(x,y)\bigr]
\]
where $H_A$ is \FOtwo-definable.
Here the quantifier binds $y$ only in the second conjunct.
\begin{restatable}[Exponential lower bound from \FOthree\ to \FOtwo\proofref{sec:FOtwo}]{theorem}{sentencebounds}\label{cor:sentences}
For every $k\ge1$,
every \FOtwo\ formula $\varphi$ such that $\varphi\equiv G_k$ has at least $2^k$ quantifiers.

%The bound holds even if $\varphi$ and $G_k$ agree only on finite structures.
\end{restatable}
The displayed expansion uses exactly $2^k$ quantifiers, so the bound
is tight. To our knowledge, this is the first exponential succinctness
gap from \FOthree\ to \FOtwo\ over general finite structures. The gap
also holds for sentences. On pure finite linear orders, by contrast,
\FOthree\ sentences admit polynomial-size \FOtwo\ translations
\citep{grohe-schweikardt2005}, while Vinall-Smeeth's exponential
separations over general structures start at three target variables
\citep{vinall-smeeth2024} (see \S\ref{sec:relwork}).

The same reference structure $\mathcal M_0$ and folds from \S\ref{sec:overview} suffice.
Indeed, regard each \FOtwo\ subformula $\psi$ as a binary relation
$X_\psi$ of assignments to $(x,y)$, with any unused coordinate treated
as a dummy. Quantifying a variable is then a special case of relational
composition:
\[
  X_{\exists y\,\psi}=X_\psi;\one,
  \qquad
  X_{\exists x\,\psi}=\one;X_\psi.
\]
Thus their assignment relations remain color-determined on $\mathcal M_0$,
just like the CoR gate values in \S\ref{sec:overview}.
Since one operand of each composition is $\one$, each quantifier can
notice at most one fold (\S\ref{ov:twocolors}), rather than the two
possible for general composition. Atomic formulas and Boolean
connectives notice none; universal quantifiers are handled by
complement. The counting inequality~\eqref{eq:fold-counting} therefore
requires at least $2^k$ quantifiers. We also obtain a lower bound of $k2^k$ relation-atom occurrences and
optimal tree size $\Theta(k2^k)$, for both queries and sentences
(\S\ref{sec:sentences}).

For \CoR\ itself, we also determine the optimal tree size of $G_k$
as $(2+o(1))k2^k$, whereas its optimal circuit size is $\Theta(2^k)$
(Theorem~\ref{thm:tree}).

\paragraph{\texorpdfstring{{\boldmath\CoR\ to \FOtwo}}{CoR to FO2}.}
An exponential succinctness gap also occurs from \CoR\ to \FOtwo.
We need a different query, since $G_k$ already requires exponentially
many \CoR\ compositions. Put $\mathscr P_i^1\coloneqq\id\cap P_i$
and $\mathscr P_i^0\coloneqq\id\cap\overline{P_i}$, and let
$\vartheta_A(u)\coloneqq\bigwedge_{i\in A}P_i(u,u)\land
\bigwedge_{i\notin A}\neg P_i(u,u)$ for $A\subseteq[k]$.
Consider the following linear-size \CoR\ term and its \FOtwo\ expansion:
\[
\defn{\CoR}:
 U_k \coloneqq R;\bigcap_{i=1}^k
 \left[(\mathscr P_i^0;\one;\mathscr P_i^0)
 \cup(\mathscr P_i^1;\one;\mathscr P_i^1)\right]
 \longrightarrow
\defn{\FOtwo}:
 \bigvee_{A\subseteq[k]}
 \left[\vartheta_A(y)\land\exists y\,
    (R(x,y)\land\vartheta_A(y))\right]
\]
Thus $U_k$ asks whether $x$ has an $R$-successor with the same
$k$-bit loop pattern as $y$.
\begin{restatable}[Exponential lower bound from \CoR\ to \FOtwo\proofref{sec:reverse}]{proposition}{reverse}\label{prop:reverse}
For every $k\ge1$, every \FOtwo\ formula
$\varphi$ such that $\varphi\equiv U_k$ has at least
$2^k$ quantifiers.
% The bound holds even if $\varphi$ and $U_k$ agree only on finite structures.
\end{restatable}

The displayed expansion uses exactly $2^k$ quantifiers, so the bound
is tight.
%
For the lower bound, take a new reference structure $\mathcal W_0$
with several vertices of each loop pattern and a hub $h$ having
$R$-successors of every pattern (Figure~\ref{fig:reverse-fold}).
The fold of $A\subseteq[k]$ deletes all $R$-edges from $h$ to
vertices of pattern $A$.
Thus $U_k(h,y_A)$ changes from true to false
for a vertex $y_A$ of pattern $A$, even when $(h,y_A)$ is unchanged.

\begin{figure}[!htbp]
\centering
\begin{tikzpicture}[
  >=Stealth,
  every node/.style={font=\small},
  patternclass/.style={draw=black!55, fill=black!3,
      rounded corners=2pt, minimum width=2.9cm, minimum height=.78cm},
  member/.style={circle, fill=black, inner sep=1.55pt},
  testmember/.style={circle, fill=black!50, inner sep=1.55pt},
  redge/.style={->, semithick, shorten >=2pt},
]
\begin{scope}
  \node[font=\bfseries\small] at (0.5, 1.0) {Reference $\mathcal W_0$};
  \node[draw,circle,inner sep=2pt] (h0) at (.20,0) {$h$};
  \node[patternclass] at (4.10,.69) {};
  \node[font=\scriptsize,anchor=south] at (4.10,1.10) {pattern $A$};
  \node[member,label={[font=\scriptsize]right:$z_A$}] (za0) at (3.54,.64) {};
  \node[testmember,label={[font=\scriptsize]right:$y_A$}] at (4.66,.64) {};
  \node[patternclass] at (4.10,-.78) {};
  \node[font=\scriptsize,anchor=south] at (4.10,-.37) {pattern $B\ne A$};
  \node[member,label={[font=\scriptsize]right:$z_B$}] (zb0) at (3.60,-.84) {};
  \draw[redge] (h0) -- (za0);
  \draw[redge] (h0) -- (zb0);
  \node[font=\scriptsize] at (2.78,-1.62) {$U_k(h,y_A)=1$};
\end{scope}
\begin{scope}[xshift=7.1cm]
  \node[font=\bfseries\small] at  (0.5, 1.0) 
  {Folded $A$: $\mathcal W_A$};
  \node[draw,circle,inner sep=2pt] (h1) at (.20,0) {$h$};
  \node[patternclass] at (4.10,.69) {};
  \node[font=\scriptsize,anchor=south] at (4.10,1.10) {pattern $A$};
  \node[member,label={[font=\scriptsize]right:$z_A$}] (za1) at (3.54,.64) {};
  \node[testmember,label={[font=\scriptsize]right:$y_A$}] at (4.66,.64) {};
  \node[patternclass] at (4.10,-.78) {};
  \node[font=\scriptsize,anchor=south] at (4.10,-.37) {pattern $B\ne A$};
  \node[member,label={[font=\scriptsize]right:$z_B$}] (zb1) at (3.60,-.84) {};
  \draw[redge,dashed,black!35] (h1) -- (za1);
  \node[font=\large, text=black!70] at (2.25,.39) {$\times$};
  \draw[redge] (h1) -- (zb1);
  \node[font=\scriptsize] at (2.78,-1.62) {$U_k(h,y_A)=0$};
\end{scope}
\end{tikzpicture}
\caption{The fold of $A$ removes all $R$-edges from $h$ to vertices of loop pattern $A$ (schematic).}
\label{fig:reverse-fold}
\end{figure}

The original folds cannot give this separation: on their reference
structure, every composition notices at most two folds, ruling out a
short \CoR\ term that distinguishes all those tests. In the new
structure, the outer composition of $U_k$ can notice all $2^k$ folds,
whereas each \FOtwo\ quantifier, viewed as a projection, still notices
at most one. The latter bound uses redundant representatives of each
atomic pair type. Applying the same counting argument~\eqref{eq:fold-counting}
to the new tests proves the lower bound (\S\ref{sec:reverse}).

\paragraph{Arity elimination in matrix query languages.}
Brijder~et~al.\ compared two formalisms for semiring-valued queries:
ARA$(3)$, which allows intermediate tables with up to three columns,
and MATLANG, which computes with matrices, vectors, and scalars.
They showed that the two have the same expressive power when inputs
and outputs have at most two columns~\cite{brijder2022}.
Their translation from ARA$(3)$ to MATLANG can grow exponentially,
but whether this blowup is unavoidable was left open~\citep[p.~9]{brijder2019}.
We establish an exponential lower bound using a semiring-valued
analogue of our hard query $G_k$.

Let $\mathfrak S=(S,\oplus,\otimes,0_S,1_S)$ be a commutative semiring,
with input relations $P_i,Q_i:V^2\to S$. Define
\[
  \mathcal G_k(x,y)\coloneqq
  \bigoplus_{z\in V}\ \bigotimes_{i=1}^k
  \Bigl[(P_i(x,y)\otimes Q_i(x,z))
  \oplus(Q_i(x,y)\otimes P_i(x,z))\Bigr].
\]
This query has an ARA$(3)$ expression of size $O(k)$: all $k$
conditions are combined over $(x,y,z)$ before $z$ is aggregated away.
Over Boolean annotations ($S=\mathbb B$, with $\oplus=\lor$ and
$\otimes=\land$), $\mathcal G_k$ specializes to $G_k$.

In MATLANG, the usual semiring matrix--vector (mv) and matrix--matrix
(mm) products take the form
\[
  (A \cdot v)(x) \coloneqq \smash{\bigoplus_{z\in V}} A(x,z)\otimes v(z)
  \quad\text{(mv)},\qquad
  (A \cdot B)(x,y) \coloneqq \smash{\bigoplus_{z\in V}} A(x,z)\otimes B(z,y)
  \quad\text{(mm)}.
\]

Both eliminate their aggregated index immediately.
The DNF of \S\ref{sec:main-results} also gives an explicit MATLANG
representation of $\mathcal G_k$. Here $H_A$ is obtained from the expression in \S\ref{sec:main-results}
by replacing intersections with pointwise $\otimes$, and
$\mathbf 1_S$ is the all-$1_S$ column vector. Then
\[
  \mathcal G_k(x,y)=\smash{\bigoplus_{A\subseteq[k]}}
       H_A(x,y)\otimes (H_{A^c}\mathbf \cdot \mathbf{1}_S)(x).
\]

This uses $2^k$ mv products, together with row broadcasting and
entrywise operations; more general computations may also use mm
products and shared intermediate results.
For a circuit $C$, let $\mvcount C$ count matrix--vector and
vector--matrix aggregations, and $\mmcount C$ count matrix--matrix
aggregations, with shared gates counted once. All other permitted
operations are free.

For a concrete non-Boolean case, consider the natural and integer
min--plus semirings
$\mathsf{Trop}_D \coloneqq (D\cup\{+\infty\},\min,+,+\infty,0)$,
where $D\in\{\mathbb N,\mathbb Z\}$.

\begin{restatable}[Tropical arity-elimination overhead\proofref{sec:matrix}]{theorem}{tropicalmatrixbounds}\label{thm:tropical-matrix}
For every $k\ge1$ and $\mathfrak S\in\{\mathsf{Trop}_{\mathbb N},\mathsf{Trop}_{\mathbb Z}\}$,
the query $\mathcal G_k$ has an ARA$(3)$ expression of size $O(k)$,
but every typed matrix circuit $C$ over $\mathfrak S$ computing it
on all finite inputs satisfies $\mvcount C+3\mmcount C\ge2^k$.
\end{restatable}
The lower bound holds even with arbitrary uniform entrywise functions
and aggregation kernels. Its proof uses the fold-counting principle
\eqref{eq:fold-counting} (\S\ref{sec:matrix}).

The argument extends uniformly beyond tropical matrices.
Examples include the Boolean semiring $(\{0,1\},\lor,\land)$
and the max--min semiring $([0,1],\max,\min)$.
It also covers integer max--plus
$(\mathbb Z\cup\{-\infty\},\max,+)$ and, for each fixed $b\ge1$,
finite-set annotations $(2^{[b]},\cup,\cap)$ and $b$-dimensional
min--plus $\mathsf{Trop}_{\mathbb Z}^{b}$.
One property covers all these examples: \emph{finite join breadth},
meaning that every finite join is determined by a bounded number
of its summands. It bounds how many folds an aggregation can notice,
giving exponential lower bounds for every fixed nontrivial
commutative idempotent semiring of finite join breadth
(Theorem~\ref{cor:matrix}). Over Boolean matrices, we also obtain
the sharper bound $\mvcount C+2\mmcount C\ge2^k$.
Positive computations over nonidempotent semirings such as
$(\mathbb N,+,\times)$ are handled separately by Booleanization
(\S\ref{sec:positive}).

\paragraph{Robustness: closure and randomization.}

The lower bound persists even when \CoR\ is extended with reflexive–transitive closure 
$X^* \coloneqq \bigcup_{m\ge 0}X^m$,
which can be evaluated in $O(n^\omega)$ time~\citep{fischer-meyer1971,munro1971,furman1970}.
Each closure gate notices at most one fold.


\begin{restatable}[Closure does not help\proofref{sec:closure}]{proposition}{closure}\label{prop:star}
Let $C$ be a \CoR\ circuit with closure gates $(-)^*$
and let $g_*(C)$ count its distinct closure gates.
For every $k \ge 1$,
if $C\equiv G_k$, then $2\gcost C+g_*(C)\ge2^k$.
\end{restatable}

The lower bound also extends to bounded-error randomized circuits, even when circuits are specialized to each sufficiently large domain size.

\begin{restatable}[Fixed sizes and bounded error\proofref{sec:query-bounds}]{corollary}{randomizedbounds}\label{cor:randomized}

For some constant $c>0$, for every $k\ge 1$, $n\ge ck4^k$, and $0\le\varepsilon<1/2$, every input-independent distribution of size-specific \CoR\ circuits whose output at each pair $(u,v)$ is correct with probability at least $1-\varepsilon$ on every $n$-element input satisfies
\[
\mathbb E[g_{;}(C)]\ge(1-2\varepsilon)2^{k-1}.
\]
\end{restatable}